% !TEX root = ../../main.tex
% C6.3 — Hiện thực giao diện người dùng (2026-09-29: 11 ảnh chụp màn hình mới H10a–k do SV chụp, images/ -> figures/*-cat.png bằng scratchpad crop_shots2.py; bản cũ ở report/archive/)
% Khớp frontend/src: routes/guards.jsx (RequireAuth, RequireRole -> trang chủ của vai trò), services/apiService.js (X-CSRF-Token cho phương thức
% không an toàn, CSRF chỉ giữ trong bộ nhớ, lỗi 5xx kèm "Mã tra cứu" = 8 ký tự đầu request id, 401 -> xóa phiên),
% store/AppStoreProvider.jsx (làm mới phiên chỉ khi có thao tác thật, chia sẻ giữa các tab), components/person/PersonCrop.jsx (khung 3:5,
% ?aspect=0.60&mark=1, "Không còn ảnh"), backend services/track_imagery.py (nới khung theo tỉ lệ bằng điểm ảnh thật, làm tối 0.6, viền;
% khung toàn cảnh vẽ khung đỏ; cạnh dài tối đa 1920; Cache-Control private, no-store; kiểm tra quyền như kết quả).
\section{Hiện thực giao diện người dùng}
\label{sec:6.3}

\subsection{Điều hướng và giao tiếp với máy chủ}
\label{sec:6.3.1}
\label{sec:6.3.2}

Ứng dụng web là một ứng dụng một trang xây dựng bằng React, với việc chuyển màn hình do React Router đảm nhận. Các đường dẫn được chia thành ba nhóm theo vai trò: nếu chưa đăng nhập, người dùng được chuyển tới màn hình đăng nhập. Nếu mở một đường dẫn không thuộc vai trò của mình, người dùng được chuyển về trang chủ của vai trò đó. Việc chặn này chỉ nhằm tránh hiển thị màn hình không dùng được, còn quyền thực sự luôn được máy chủ kiểm tra ở từng yêu cầu.

Mọi yêu cầu tới máy chủ đi qua một lớp gọi API chung. Lớp này gắn mã CSRF vào mọi yêu cầu thay đổi dữ liệu (Mục~\ref{sec:5.5.2}). Mã CSRF chỉ được giữ trong bộ nhớ của trang, không lưu vào bộ nhớ trình duyệt. Phản hồi lỗi được chuyển thành thông báo tiếng Việt. Với lỗi phía máy chủ, thông báo kèm một \emph{mã tra cứu} gồm tám ký tự đầu của mã yêu cầu để Quản trị viên tìm đúng dòng nhật ký (Mục~\ref{sec:6.4}). Khi bất kỳ yêu cầu nào nhận mã 401, ứng dụng xóa trạng thái đăng nhập và đưa người dùng về màn hình đăng nhập. Phiên chỉ được làm mới khi người dùng thực sự thao tác. Các yêu cầu tự động chạy nền không giữ phiên sống, và thời điểm thao tác gần nhất được chia sẻ giữa các thẻ trình duyệt.

\subsection{Hiển thị ảnh người trong kết quả}
\label{sec:6.3.3}

Ảnh người trong danh sách kết quả được máy chủ dựng từ khung hình toàn cảnh và khung bao mỗi khi được yêu cầu (Mục~\ref{sec:5.3.3}). Thẻ kết quả trên giao diện có tỉ lệ 3:5; vùng cắt được nới rộng từ khung bao bằng phần cảnh thật xung quanh cho tới khi đạt tỉ lệ này, chỉ nới ra chứ không cắt vào người. Vì vùng cắt rộng có thể chứa cả người khác, phần xung quanh được giảm độ sáng còn 60\% và người được tìm thấy được viền bằng màu nhấn. Khung hình toàn cảnh trong hộp thoại chi tiết được vẽ khung bao màu đỏ quanh người. Ảnh chỉ được trả khi người yêu cầu có quyền xem kết quả tương ứng, kèm chỉ dẫn không lưu vào bộ nhớ đệm; nếu khung hình không còn trong kho, thẻ kết quả hiển thị biểu tượng \emph{Không còn ảnh}.

\subsection{Các màn hình chính}
\label{sec:6.3.4}

Các hình dưới đây là những màn hình chính của ứng dụng theo từng vai trò, chụp trên dữ liệu trình diễn (Mục~\ref{sec:7.2}).

\paragraph{Giám sát viên.} Màn hình \emph{Tìm kiếm người} (Hình~\ref{fig:ui-search-image}) gồm biểu mẫu truy vấn bên trái, với ảnh truy vấn, các camera trong khu vực, khoảng ngày và số kết quả, và danh sách kết quả xếp hạng bên phải. Mỗi kết quả hiển thị ảnh người trong ngữ cảnh, camera, giờ xuất hiện và điểm phù hợp; khu vực và ngày chung chỉ hiển thị một lần ở đầu danh sách, và danh sách có thể lọc lại theo thời gian, theo camera hay sắp xếp theo thời gian mà không gửi truy vấn mới. Khi tìm theo thuộc tính (Hình~\ref{fig:ui-search-attr}), câu mô tả tiếng Anh được sinh từ các thuộc tính đã chọn hiển thị ngay trong biểu mẫu, để người dùng biết hệ thống thực sự tìm theo mô tả nào. Chọn một kết quả mở hộp thoại chi tiết với khung hình toàn cảnh (Hình~\ref{fig:ui-result}), từ đó kết quả được lưu vào vụ việc mới, với người phụ trách tự động là tài khoản đang đăng nhập và tùy chọn đánh dấu hoàn thành ngay (Hình~\ref{fig:ui-create-case}). Màn hình \emph{Vụ việc của tôi} (Hình~\ref{fig:ui-cases}) gồm danh sách vụ việc theo trạng thái và chi tiết vụ việc đang chọn; vụ việc đã hoàn thành bị khóa và chỉ có thể mở lại.

\begin{figure}[!htbp]
\centering
\includegraphics[width=\textwidth]{H10a-tim-kiem-anh-cat.png}
\caption{Màn hình tìm kiếm người bằng ảnh}
\label{fig:ui-search-image}
\end{figure}

\begin{figure}[!htbp]
\centering
\includegraphics[width=\textwidth]{H10b-tim-kiem-thuoc-tinh-cat.png}
\caption{Màn hình tìm kiếm người theo thuộc tính}
\label{fig:ui-search-attr}
\end{figure}

\begin{figure}[!htbp]
\centering
\begin{subfigure}[b]{0.61\textwidth}
    \centering
    \includegraphics[width=\linewidth]{H10c-chi-tiet-ket-qua-cat.png}
    \caption{Chi tiết một kết quả}
    \label{fig:ui-result}
\end{subfigure}\hfill
\begin{subfigure}[b]{0.36\textwidth}
    \centering
    \includegraphics[width=\linewidth]{H10d-tao-vu-viec-cat.png}
    \caption{Tạo vụ việc mới từ kết quả}
    \label{fig:ui-create-case}
\end{subfigure}
\caption{Hộp thoại chi tiết kết quả và hộp thoại tạo vụ việc}
\label{fig:ui-result-case}
\end{figure}

\begin{figure}[!htbp]
\centering
\includegraphics[width=\textwidth]{H10e-vu-viec-cua-toi-cat.png}
\caption{Màn hình quản lý vụ việc của Giám sát viên}
\label{fig:ui-cases}
\end{figure}

\paragraph{Quản lý.} Màn hình \emph{Tổng quan} (Hình~\ref{fig:ui-overview}) hiển thị số vụ việc theo trạng thái, số kết quả đã lưu và danh sách vụ việc gần đây; chọn một vụ việc mở màn hình \emph{Hồ sơ vụ việc}, có cùng bố cục như Hình~\ref{fig:ui-cases} nhưng ở chế độ chỉ đọc.

\begin{figure}[!htbp]
\centering
\includegraphics[width=\textwidth]{H10f-tong-quan-cat.png}
\caption{Màn hình tổng quan của Quản lý}
\label{fig:ui-overview}
\end{figure}

\paragraph{Quản trị viên.} Màn hình \emph{Camera} (Hình~\ref{fig:ui-cameras}) liệt kê các camera theo khu vực cùng trạng thái kết nối, trong đó camera chỉ xử lý tệp video được ghi là \emph{Video tải lên}, trạng thái xử lý AI và các thao tác kiểm tra kết nối, sửa, loại khỏi vận hành. Màn hình \emph{Mô hình AI} (Hình~\ref{fig:ui-models}) cho chọn Detector và Tracker; mô hình không đủ điều kiện được đánh dấu \emph{Không khả dụng}, còn bộ mã hóa được hiển thị là thành phần cố định (Mục~\ref{sec:6.2.1}). Màn hình \emph{Xử lý video} (Hình~\ref{fig:ui-video}) dùng để tải video lên cho một camera với chế độ lấy mẫu định sẵn và theo dõi tiến độ từng công việc theo số khung hình và số track đã sẵn sàng. Hai màn hình giám sát là \emph{Trạng thái hệ thống} (Hình~\ref{fig:ui-status}, Mục~\ref{sec:6.4}) và \emph{Kiểm tra AI} (Hình~\ref{fig:ui-diagnostics}, Mục~\ref{sec:6.2.5}).

\begin{figure}[!htbp]
\centering
\includegraphics[width=\textwidth]{H10g-camera-cat.png}
\caption{Màn hình quản lý camera}
\label{fig:ui-cameras}
\end{figure}

\begin{figure}[!htbp]
\centering
\includegraphics[width=\textwidth]{H10h-mo-hinh-ai-cat.png}
\caption{Màn hình cấu hình mô hình AI}
\label{fig:ui-models}
\end{figure}

\begin{figure}[!htbp]
\centering
\includegraphics[width=\textwidth]{H10i-xu-ly-video-cat.png}
\caption{Màn hình xử lý video tải lên}
\label{fig:ui-video}
\end{figure}

\begin{figure}[!htbp]
\centering
\includegraphics[width=\textwidth]{H10j-trang-thai-he-thong-cat.png}
\caption{Màn hình trạng thái hệ thống}
\label{fig:ui-status}
\end{figure}

\begin{figure}[!htbp]
\centering
\includegraphics[width=\textwidth]{H10k-kiem-tra-ai-cat.png}
\caption{Màn hình kiểm tra hoạt động của AI}
\label{fig:ui-diagnostics}
\end{figure}
